% ID: SEL-D0033-OA-E03
% ESTADO: cotejada. Segunda pasada visual y borrador compilado revisados.
% Fuente: Selectividad/matematicas_aplicadas/titular_septiembre/Examen-Matemáticas Aplicadas a las Ciencias Sociales.pdf
% SHA-256: 0f95b6aefa78f1a13e2115b281338bd807177671c303c380180957fe03f0a605
% Página: 1; opción A; ejercicio original 3.

En una localidad, el 25\% de los habitantes asiste periódicamente a la consulta del dentista, el 10\% se hace una analítica y el 8\% hace ambas cosas.
\begin{enumerate}[label=\alph*)]
\item \textbf{(0.5 puntos)} Razone si los sucesos ``Asistir a la consulta del dentista'' y ``Hacerse una analítica'' son independientes.
\item \textbf{(1 punto)} \textquestiondown{}Qué porcentaje de habitantes ni se hace una analítica ni va al dentista?
\item \textbf{(1 punto)} Si elegimos un habitante al azar de esa localidad de entre los que no van al dentista, \textquestiondown{}cuál es la probabilidad de que se haga una analítica?
\end{enumerate}
